\section{The complete exterior and its energy}\label{n3:limit:section}

We now obtain the three geometric outputs stated at the beginning:
nonnegative scalar curvature, control of every cut, and an arbitrarily
small upper energy error. Their combination will also yield the
numerical comparison. Fix the prepared initial data of
Definition~\ref{n3:prepared:data}, and the prepared exterior
\(\Omega\) of Proposition~\ref{n3:domains:prepared-domain}.
Its boundary \(B\) separates the original inner obstacle from the end.
Let \(A_*\) be the enclosing-area infimum of that original obstacle,
measured in the reduced metric \(g\).
In particular every cut enclosing \(B\) has \(g\)-area at least \(A_*\).

We use the notation of Section~\ref{n3:deformation:section} and take the
physical endpoint of the homotopy:
\[
 F=h+C,\qquad h=\tau f,\qquad
 \Div V=u\Xi,\qquad
 \Xi=\delta_0\mathcal T+\rho+\delta_0\tau a\sigma-m_0(t)\mathcal P.
\]
Here \(0<\delta_0<1\), \(\ell=e^{-\epsilon N}\),
\(\tau=\ell^{3/2}\), and all constants in the prepared data are fixed
before \(N\) is taken large. Constants marked \(C_N\) in this section
are bounded by \(C(1+N)^C\), with fixed exponents. Other constants
explicitly allowed to depend on fixed \(N\) need not have such a bound.

\subsection{The limiting solution and its end}

\begin{lemma}[End control]\label{n3:limit:end}
Fix \(\epsilon>0\) and a sufficiently large
\(N\ge\max\{4,\lfloor2/\epsilon\rfloor+1\}\). The truncation solutions of
Theorem~\ref{n3:existence:deformation} have a subsequence converging
smoothly on compact subsets of \(\overline\Omega\) to a solution of
the physical system with \(t\ge-\epsilon\).
On the asymptotically flat end, for some \(c_N>0\) and fixed-\(N\)
constants \(C_j(N)\),
\[
 |\nabla^j f|\le C_j(N)e^{-c_N r}
 \quad\text{for each fixed }j,\qquad
 Z,t=O_2(r^{-1}).
\]
The differences \(t-Z\), \(\bar g-g\), and their coordinate derivatives
of any fixed order decay exponentially. The vector field \(V\)
satisfies
\begin{equation}\label{n3:limit:V-asymptote}
                    V=4\nabla_g t+O(r^{-3})
\end{equation}
and has a finite outward flux limit.
\end{lemma}

\begin{proof}
The compact convergence, including at the fixed inner boundary, follows
from Theorem~\ref{n3:existence:deformation} and its uniform fixed-\(N\)
classical estimates. All equations and inner boundary conditions pass
to this limit. To preserve the end normalization, we prove the decay
estimates on the truncations before passing to infinity.

Outside a fixed large sphere, \(K=C=0\), and the first equation is
\[
 \tr_{A_\chi}\Hess_g f=\frac{\tau\sqrt D}{l}\,f.
\]
At fixed \(N\) the preceding estimates make the left side uniformly
elliptic, and its positive height coefficient is bounded below by a
positive constant depending on \(N\). It is bounded above as well.
For sufficiently small \(c_N>0\), the functions
\(\exp[-c_N(r-r_0)]\) are positive supersolutions of the corresponding
linear equation outside a sufficiently large \(r_0\): their radial
second derivative has size \(c_N^2\), while their tangential second
derivatives have size \(c_N/r\), and the fixed positive height term
dominates both. Multiplying by a constant bounds the solution on
\(S_{r_0}\); the zero outer Dirichlet value is also bounded.
The maximum principle applied to both signs of \(f\) gives the
uniform exponential estimate. The fixed-\(N\) unit-patch estimates,
followed by the local homogeneous linear estimates for this equation,
give the same decay for derivatives. The corresponding boundary
estimates apply on the large outer spheres. Iterating the smooth
equations gives each required finite derivative order.

The relation \(Z=t+\tfrac18\log(1+l^2|\nabla f|^2)\) now shows that
\(t-Z\) is exponentially small, with derivatives. The graph errors
have the same property. On the end the second equation becomes
\begin{equation}\label{n3:limit:Z-end}
 \Delta_gZ+b_N(Z)|dZ|_g^2=O(r^{-4}),\qquad
 b_N(s)=p_L(s)-\delta_0\bigl(p_L(s)-1\bigr).
\end{equation}
Indeed \(A_\chi=I\) and \(u=L(Z)\) up to exponentially small errors,
whereas
\(\mathcal T=4(p_L(Z)-1)|dZ|^2\) up to such errors.
The \(\rho_0\) penalty weight is \(O(r^{-4})\); the term with
\(\rho_1\) is multiplied by the exponentially small \(v\).
All coefficients in this statement are bounded on the fixed-\(N\)
solution range. Differentiably controlled exponential errors can
absorb any fixed polynomial power of \(r\).

Define an increasing smooth function on that range by
\[
 \Phi_N(0)=0,\qquad
 \Phi_N'(s)=\exp\!\left(\int_0^s b_N(z)\,\mathrm dz\right).
\]
Its derivative is bounded above and below by positive fixed-\(N\)
constants. The chain rule in \eqref{n3:limit:Z-end} gives
\[
                         \Delta_g\Phi_N(Z)=O(r^{-4}).
\]
The function is bounded and vanishes at the outer cutoff. To obtain
a bound independent of that cutoff, choose \(0<\eta<1\).
For large \(r\), a positive multiple of
\(r^{-1}-r^{-1-\eta}\) satisfies
\[
 -\Delta_g(r^{-1}-r^{-1-\eta})\ge c r^{-3-\eta}.
\]
Here the Euclidean leading term is
\(\eta(1+\eta)r^{-3-\eta}\), and the metric error is \(O(r^{-4})\).
It dominates the right side after increasing the multiple and the
fixed inner radius. Comparison for both signs yields
\(|\Phi_N(Z)|\le C'(N) r^{-1}\), where \(C'(N)\) is a fixed-\(N\)
constant with no polynomial claim. Consequently \(Z=O(r^{-1})\).

At sufficiently large radius this estimate and the exponential graph
error give \(t>-\epsilon+1/N\); hence the floor penalty vanishes there.
Rescale each annulus of radii comparable to \(r\) to unit size.
The equation for \(\Phi_N(Z)\) has a bounded rescaled source and
uniformly controlled metric coefficients. Local \(W^{2,p}\)
estimates with \(p>3\) first give the gradient decay. Its source is
then a smooth weight times a smooth function of the bounded variables,
plus the already controlled exponential errors. Schauder estimates,
or a further Sobolev bootstrap followed by them, give
\(Z=O_2(r^{-1})\). Inverting \(\Phi_N\) preserves these bounds, and
the same holds for \(t\).

Since \(L(0)=1\), we have \(u=1+O(r^{-1})\) and
\(A_\chi=I+O(e^{-c_N r})\). With \(K=0\) on the end, the definition
of \(V\) gives \eqref{n3:limit:V-asymptote}.
Moreover \(\Div V=u\Xi\) is integrable there:
\(\mathcal T=O(r^{-4})\), \(\rho=O(r^{-4})\), the penalty has vanished,
and \(\tau a\sigma\) decays exponentially. The divergence theorem on
annuli proves the existence of the outward flux limit.
\end{proof}

\begin{lemma}[Curvature, completeness, and ADM charge]\label{n3:limit:geometry}
For the solution in Lemma~\ref{n3:limit:end},
\[
                         \hat g=e^{4t}(g+l^2df^2)
\]
is a smooth complete metric on \(\overline\Omega\), with nonnegative
scalar curvature and strictly negative mean curvature on \(B\) for
the normal pointing into \(\Omega\).
Its end is asymptotically flat with
\[
 \hat g-\delta=O_2(r^{-1}),\qquad
 R_{\hat g}=O(r^{-3-\delta'}),\qquad
 \delta'=\min\{\delta,1\}>0,
\]
where the reduced end has \(R_g=O(r^{-3-\delta})\).
If
\[
                 \mathfrak F_N=\lim_{r\to\infty}
                    \int_{S_r}V_\nu\dd A_g,
\]
then
\begin{equation}\label{n3:limit:adm}
                         E_{\hat g}=E-\frac{\mathfrak F_N}{8\pi}.
\end{equation}
\end{lemma}

\begin{proof}
At the physical endpoint the curvature identity gives
\begin{align*}
 \tfrac12e^{4t}R_{\hat g}
 ={}&8\pi(\mu+J(w))+(1-\delta_0)\mathcal T
       +(1-\delta_0)\tau a\sigma\\
    &+w\cdot\nabla C-F\tr K-\rho+m_0(t)\mathcal P.
\end{align*}
The height bounds on the compact supports of \(K\) and \(C\) give
\(h\in[b_-,b_+]\). Since \(|w|\le1\), the terms
\(|F\tr K|+|\nabla C|+\rho\) are dominated by
\(8\pi(\mu-|J|)\), by the choices in
Proposition~\ref{n3:domains:prepared-domain} and the choice of \(\rho\).
Outside those supports the same conclusion uses
\(8\pi(\mu-|J|)-\rho\ge0\).
All remaining terms are nonnegative. This proves
\(R_{\hat g}\ge0\).

The boundary identity and the prescribed flux give
\(H+4\partial_\nu t=-N_B<0\). Since \(f\) is constant on each face,
this is exactly the mean-curvature sign for \(\hat g\).
Also \(\hat g\ge e^{-4\epsilon}g\), so any escaping curve has infinite
length; the compact boundary is included. Smoothness and completeness
therefore follow.

Lemma~\ref{n3:limit:end} gives the stated metric decay. The prepared-data hypothesis supplies \(R_g=O(r^{-3-\delta})\) on the end. Since graph errors are exponentially small, the
conformal formula gives
\[
 R_{\hat g}
 =e^{-4t}\bigl(R_g-8\Delta_gt-8|dt|_g^2\bigr)
      +O(e^{-c_N r}).
\]
Equation~\eqref{n3:limit:Z-end}, the vanishing far-end penalty, and the
gradient estimate give \(\Delta_gt=O(r^{-4})\).
This proves the asserted scalar-curvature decay, including the
pointwise falloff required by the AF convention of
\cite[Definition~21 and Theorem~19]{Bray2001}.

Exponential graph errors contribute no ADM flux. The conformal change
and \(t=O_2(r^{-1})\) give
\[
 E_{\hat g}
 =E-\frac1{2\pi}\lim_{r\to\infty}
                 \int_{S_r}\partial_\nu t\dd A_g.
\]
Replacing the metric normal and area form by their Euclidean versions
costs \(o(1)\), and the \(O(r^{-3})\) error in
\eqref{n3:limit:V-asymptote} has vanishing sphere integral.
This proves \eqref{n3:limit:adm}.
\end{proof}

\subsection{The boundary and floor losses}

The mass formula reduces the remaining estimate to a lower bound for
the flux. We absorb the inner boundary term into the positive bulk
source and then estimate the penalty near the floor. Both steps use
only polynomial constants.

\begin{lemma}[Flux lower bound]\label{n3:limit:flux}
For fixed \(\epsilon>0\), the physical solutions satisfy
\[
 \mathfrak F_N\ge-\varepsilon_N,\qquad
 0\le\varepsilon_N\le
      C(1+N)^C\bigl(\ell+\ell^2/\tau\bigr)\longrightarrow0.
\]
All constants in this estimate are independent of the fixed-\(N\)
classical regularity constants.
\end{lemma}

\begin{proof}
The normal \(\nu\) on \(B\) points into \(\Omega\), so the divergence
theorem has the sign
\begin{equation}\label{n3:limit:flux-split}
 \mathfrak F_N
     =\int_\Omega u\Xi\dd V_g+\int_BV_\nu\dd A_g.
\end{equation}
Each integral is finite for a fixed \(N\), by
Lemma~\ref{n3:limit:end} and compact smoothness.

At a black face, \(F-P_B=H\) and \(w_\nu=a\).
At a white face, \(F-P_B=-H\) and \(w_\nu=-a\).
Thus at either type of face,
\[
                 V_\nu/u=-(1-a)H-N_Bd\ge-Cd.
\]
The signed boundary gradient bounds give
\(\sigma\ge c/\tau\), so
\[
 \sqrt d\le\frac1{l\sigma}\le C\tau/\ell,\qquad
 \int_BV_\nu\dd A_g\ge
             -C\frac{\tau}{\ell}\int_BL\dd A_g.
\]
Use the polynomial collar trace estimate from
Section~\ref{n3:apriori:section} with the constant test function:
\[
 \int_B L\dd A_g
    \le C(1+N)^C
       \int_{\mathcal C}u(1+\sqrt{\mathcal T})\dd V_g,
\]
where \(\mathcal C\) is a fixed compact union of collars.
Its polynomial constant is independent of the upper \(Z\) bound.
Since \(\rho\) has a fixed positive minimum on \(\mathcal C\),
\[
 \int_{\mathcal C}u(1+\sqrt{\mathcal T})\dd V_g
 \le C\int_\Omega u(\delta_0\mathcal T+\rho)\dd V_g.
\]
The factor \(C(1+N)^C\tau/\ell
=C(1+N)^C\ell^{1/2}\) tends to zero.
For large \(N\), the entire negative boundary term in
\eqref{n3:limit:flux-split} is therefore at most one quarter of
\(\int_\Omega u(\delta_0\mathcal T+\rho)\dd V_g\).

It remains to estimate the floor penalty. On its support,
\(-\epsilon\le t\le-\epsilon+1/N\), so \(l\) and \(L\) are comparable
to \(\ell\), with constants independent of \(N\).
Directly from their definitions,
\[
 \begin{split}
 u&=L\sqrt{1+l^2\sigma^2}
                      \le C(\ell+\ell^2\sigma),\\
 uv&=\frac{Ll^2\sigma^2}{\sqrt{1+l^2\sigma^2}}
                      \le C\ell^2\sigma,\\
 ua\sigma&=Ll\sigma^2\ge c\ell^2\sigma^2 .
 \end{split}
\]
Consequently
\[
 um_0\mathcal P
 \le C_N\bigl[\ell\rho_0+\ell^2\sigma(\rho_0+\rho_1)\bigr].
\]
Young's inequality, with the square term chosen to be a quarter of
\(\delta_0\tau ua\sigma\), yields
\[
 um_0\mathcal P
 \le \tfrac14\delta_0\tau ua\sigma+
 C_N\left[\ell\rho_0+
              \frac{\ell^2}{\tau}(\rho_0^2+\rho_1^2)\right].
\]
Squaring the constant only changes the polynomial
\(C_N\). The integrals of \(\rho_0\), \(\rho_0^2\), and \(\rho_1^2\)
are finite: on the end their orders are respectively \(r^{-4}\),
\(r^{-8}\), and \(r^{-4\gamma}\), with \(4\gamma>3\).
Integrating leaves at most
\(C(1+N)^C(\ell+\ell^2/\tau)\), after absorption.

Substitute both bounds in \eqref{n3:limit:flux-split}.
Positive fractions of the three bulk terms
\(\delta_0\mathcal T,\rho,\delta_0\tau a\sigma\) remain.
Dropping them proves the claim. Finally
\(\ell^2/\tau=\ell^{1/2}\), so the error tends to zero
exponentially against every polynomial loss.
\end{proof}

\begin{proof}[Proof of Theorem~\ref{n3:prepared:output}]
The fixed subexterior and preservation of every full cut are supplied
by Proposition~\ref{n3:domains:prepared-domain}. Fix $\epsilon$ and
take the global solutions of Lemma~\ref{n3:limit:end}, obtained by
first exhausting the end at fixed $N$. Lemma~\ref{n3:limit:geometry}
gives completeness, the two curvature signs, and the end decay. The
pointwise inequality $\widehat g\ge e^{-4\epsilon}g$ gives the
two-dimensional area inequality on every full cut, including every
component and every coincident portion. The energy formula and
Lemma~\ref{n3:limit:flux} give the asserted upper error with
$\eta_{\epsilon,N}=\varepsilon_N/(8\pi)$. This completes the
geometric construction before any Riemannian Penrose comparison.
\end{proof}

\subsection{The Riemannian horizon and the inequality}

The deformed metric has nonnegative scalar curvature, strictly negative
inner mean curvature toward the end, and the required asymptotic decay.
We replace its inner boundary by an outer area-minimizing minimal
enclosure, then compare its area with the original enclosing infimum.

\begin{theorem}[Reduced energy inequality]\label{n3:limit:energy}
The reduced data satisfy
\[
                              E\ge\sqrt{\frac{A_*}{16\pi}}.
\]
Consequently Theorem~\ref{n3:intro:exterior-inequality} holds in its
entire stated exterior class.
\end{theorem}

\begin{proof}
For each fixed sufficiently large \(N\), minimize \(\hat g\)-area
among filled cuts enclosing \(B\), using a large outer sphere as an
auxiliary barrier. The independent enclosure results of~\cite[Lemmas~2.1--2.3]{RiemannianCompanion}
apply to this smooth AF metric; we check their hypotheses and the
strict detachment here.
A fixed enclosing competitor bounds the area, and local perimeter
density estimates prevent minimizers from escaping to infinity.
The positive mean curvature of sufficiently large outer spheres
excludes contact there, as in that enclosure theorem.
For the inner barrier, let \(U_s\) be a short outward normal collar
of the filled obstacle and \(Y\) its outward unit foliation field.
Strict negativity gives \(\Div_{\hat g}Y<0\) throughout the collar.
If a perimeter minimizer \(E\) omitted a positive-volume portion
\(D=U_s\setminus E\), Gauss--Green and \(|Y|_{\hat g}\le1\) would give
\[
 P_{\hat g}(E\cup U_s)-P_{\hat g}(E)
       \le\int_D\Div_{\hat g}Y\dd V_{\hat g}<0.
\]
Thus \(E\) contains a whole inner collar; obstacle regularity upgrades
this inclusion to its regular representative. Its free boundary is
disjoint from the obstacle. The resulting cut \(\Gamma_N\) is compact,
smooth, and minimal in dimension three; it may be disconnected.
It is outer area-minimizing by its minimizing property.
Equivalently one may first use the ordinary trapped-region theorem
for \((\hat g,0)\) and then take the outer minimizing enclosure.
The boundary version of the Riemannian Penrose theorem applies to the
exterior of \(\Gamma_N\), by Lemma~\ref{n3:limit:geometry}
and \cite[Theorem~19]{Bray2001}.

For every two-plane the restriction of \(\bar g=g+l^2df^2\) dominates
that of \(g\). The floor therefore gives
\[
                  |\Gamma_N|_{\hat g}
          \ge e^{-4\epsilon}|\Gamma_N|_g
          \ge e^{-4\epsilon}A_*.
\]
The last inequality uses that \(B\), and hence \(\Gamma_N\), cuts off
the original filled obstacle. Lemmas~\ref{n3:limit:geometry}
and~\ref{n3:limit:flux} now give
\[
 E+\frac{\varepsilon_N}{8\pi}
       \ge E_{\hat g}
       \ge\sqrt{\frac{|\Gamma_N|_{\hat g}}{16\pi}}
       \ge e^{-2\epsilon}\sqrt{\frac{A_*}{16\pi}}.
\]
Let \(N\to\infty\) with \(\epsilon\) fixed, and then
\(\epsilon\downarrow0\). This proves the reduced inequality.

Finally apply the weak three-dimensional rest-preparation theorem
of~\cite[Proposition~7.13]{NumericalCompanion} to an exterior satisfying the hypotheses
of Theorem~\ref{n3:intro:exterior-inequality}.
Write $(g_j,K_j)$ for the resulting data. Their tensor fields have
compact support, their momenta vanish, and their metrics have all-order
$r^{-1}$ decay and the stipulated scalar falloff. Discard a finite
initial segment so that $E_j>0$. Strict trapping and strict DEC hold
on the compact core. The positive minimum there of
$8\pi(\mu_j-|J_j|_{g_j})\varrho^{3+\delta}$, together with the
end margin, gives the global weighted bound after decreasing $c_j$.
The full enclosing area is positive by the geometric enclosure
lemma. Thus each fixed $j$ satisfies Definition~\ref{n3:prepared:data};
no constant uniform in $j$ has been used in its deformation solve.
The reduced energies converge to its invariant ADM mass and the
full enclosing infima converge to the original enclosing infimum.
Applying the reduced bound to each member of that sequence and passing
to the limit gives Equation~\eqref{n3:intro:penrose}.
\end{proof}
